\medskip

\begin{enumerate}
\item On sait que $x \e^{-x} = \dfrac{x}{\e^x}$ et que $\displaystyle\lim_{x \to + \infty} \dfrac{x}{\e^x} = 0$, donc $\displaystyle\lim_{x \to + \infty}  f(x) = 0$ : géométriquement l'axe des abscisses est asymptote horizontale à la représentation graphique de $f$ au voisinage de plus l'infini. L'affirmation 1 est vraie.

On a $f'(x) = 5\e^{-x}  - 5x\e^{-x} = \e^{-x}(5 - 5x) = 5\e^{-x}(1 - x)$.

$f'(x) + f(x) = 5\e^{-x}(1 - x) + 5x\e^{-x}  = 5\e^{-x}$, donc $f$ vérifie l'équation différentielle $(E)$ : l'affirmation 2 est vraie.
\item La suite $\left(v_n\right)$ est minorée et majorée mais on ne sait rien de sa monotonie, rien sur sa convergence : l'affirmation 3 est fausse.

\emph{Exemple} : $u_n = - 1 - \dfrac 1n, \: v_n = (-1)^n$ et $w_n = 1 + \dfrac 1n$ : on ne peut rien dire de la suite $\left(v_n\right)$.

Si $\left(u_n\right)$ est croissante et $\left(w_n\right)$ décroissante, on a donc 

$u_0 \leqslant u_1 \leqslant \ldots u_n \leqslant v_n \leqslant w_n \leqslant w_{n - 1} \leqslant \ldots \leqslant w_1 \leqslant w_0$ : on a bien l'encadrement considéré : l'affirmation 4 est vraie.
\end{enumerate}

\bigskip
